This chapter reports the results of the study. Unless stated otherwise, every
number is measured on the held-out ISIC-2019 test split (4{,}796 images) at the
eight-view test-time-augmentation (TTA) setting, using the checkpoint selected by
the best validation BMA. All values are from a single run with seed 42.

\section{Performance Evaluation}

Table~\ref{tab:main} compares the two proposed models against the five
lightweight baselines, with every metric reported at the same TTA inference
setting so the rows are directly comparable. The lightweight hybrid reaches a BMA
of 0.6081 with only 3.98M parameters, which is the best balanced accuracy of any
model in the sub-6M-parameter budget --- ahead of MobileViT-S (0.6020), its
closest competitor, at fewer parameters and about two-thirds of the compute.
The flagship DEKAN model reaches 0.6438 BMA, the best in the study, and also
leads on every secondary metric (macro-F1 0.6059, macro-AUC 0.9271, accuracy
0.7467), at the cost of roughly four times the parameters and ten times the
compute of the lightweight model. Figure~\ref{fig:comparison} shows the same
comparison graphically.

\begin{table}[H]
\centering
\caption{Main results on the ISIC-2019 test set (with 8-view TTA). Best value in
each column is in \textbf{bold}; the two proposed models are marked.}
\label{tab:main}
\begin{tabular}{lrrrrrr}
\toprule
Model & Params (M) & GMACs & BMA & Macro-F1 & Macro-AUC & Accuracy \\
\midrule
ResNet-18            & 11.69 & 1.810 & 0.3368 & 0.2763 & 0.8217 & 0.5225 \\
MobileNetV2          & 3.50  & 0.300 & 0.4738 & 0.4701 & 0.8827 & 0.6731 \\
EfficientNet-B0      & 5.29  & 0.390 & 0.4946 & 0.4736 & 0.9037 & 0.6779 \\
MobileViT-S          & 5.60  & 1.000 & 0.6020 & 0.5354 & 0.9218 & 0.6981 \\
EfficientFormer-L1   & 12.27 & 1.300 & 0.5519 & 0.5779 & 0.9005 & 0.7333 \\
\midrule
Hybrid (ours, light) & 3.98  & 0.631 & 0.6081 & 0.5602 & 0.9229 & 0.7162 \\
DEKAN (ours, flagship) & 16.45 & 6.633 & \textbf{0.6438} & \textbf{0.6059} & \textbf{0.9271} & \textbf{0.7467} \\
\bottomrule
\end{tabular}
\end{table}

\begin{figure}[H]
\centering
\includegraphics[width=\textwidth]{02_model_comparison.pdf}
\caption{Model comparison on the ISIC-2019 test set: BMA (with and without TTA)
on the left, macro-F1 and macro-AUC on the right. The two proposed models are
highlighted.}
\label{fig:comparison}
\end{figure}

The training behaviour of the lightweight hybrid is shown in
Figure~\ref{fig:curves}. Training and validation loss decrease smoothly and the
validation BMA peaks at 0.6505 at epoch 149, with no sign of the rare-class
collapse that a poorly tuned imbalance loss would cause --- evidence that the
$\beta=0.999$ choice in the Class-Balanced Focal Loss was appropriate.

\begin{figure}[H]
\centering
\includegraphics[width=\textwidth]{01_training_curves.pdf}
\caption{Training and validation loss (left) and BMA (right) for the lightweight
hybrid model over 150 epochs. The dashed line marks the best validation BMA.}
\label{fig:curves}
\end{figure}

\section{Per-Class Analysis}

Because BMA is an average over classes, it is worth looking at where the models
succeed and fail per class. Table~\ref{tab:perclass} lists per-class recall for
two representative baselines and the two proposed models, and
Figure~\ref{fig:perclass} compares the lightweight hybrid against EfficientNet-B0
class by class.

\begin{table}[H]
\centering
\caption{Per-class recall on the ISIC-2019 test set (with TTA).}
\label{tab:perclass}
\begin{tabular}{lrrrr}
\toprule
Class & EfficientNet-B0 & MobileViT-S & Hybrid (ours) & DEKAN (ours) \\
\midrule
MEL  & 0.594 & 0.447 & 0.518 & 0.611 \\
NV   & 0.857 & 0.890 & 0.859 & 0.859 \\
BCC  & 0.546 & 0.467 & 0.618 & 0.731 \\
AK   & 0.348 & 0.452 & 0.587 & 0.368 \\
BKL  & 0.366 & 0.587 & 0.632 & 0.622 \\
DF   & 0.245 & 0.660 & 0.566 & 0.509 \\
VASC & 0.843 & 0.882 & 0.863 & 0.922 \\
SCC  & 0.157 & 0.430 & 0.223 & 0.529 \\
\bottomrule
\end{tabular}
\end{table}

The pattern is consistent with the headline numbers. The lightweight hybrid lifts
several of the hard middle classes relative to the convolutional baselines ---
BCC, AK, and BKL in particular --- which is where its global reasoning and
metadata context help most. Its weakest class is SCC (0.223), a rare class that
is easily confused with AK and BKL. DEKAN spreads its recall more evenly and
recovers SCC substantially (0.529) while also posting the strongest VASC (0.922)
and BCC (0.731), which is the main reason its balanced accuracy is highest.
Figure~\ref{fig:confusion} shows the full confusion matrix for the lightweight
hybrid; the off-diagonal mass sits mostly among the keratinocytic classes
(AK/BKL/SCC), exactly the clinically difficult group.

\begin{figure}[H]
\centering
\includegraphics[width=0.85\textwidth]{04_per_class_recall.pdf}
\caption{Per-class recall of the lightweight hybrid versus EfficientNet-B0,
sorted by the hybrid's recall.}
\label{fig:perclass}
\end{figure}

\begin{figure}[H]
\centering
\includegraphics[width=0.8\textwidth]{03_confusion_matrix.pdf}
\caption{Confusion matrix for the lightweight hybrid (normalised by true class,
with TTA).}
\label{fig:confusion}
\end{figure}

\section{Efficiency Analysis}

The point of the lightweight model is accuracy per unit of cost, so
Table~\ref{tab:efficiency} reports balanced accuracy per million parameters
alongside the raw numbers, and Figure~\ref{fig:efficiency} plots accuracy against
parameter count with compute shown as bubble size. Among all models that fit the
sub-6M budget, the lightweight hybrid has the best BMA-per-parameter (0.153),
and the CNN-stem-only ablation is even higher per-parameter (0.268) but at a much
lower absolute BMA. DEKAN sits at the opposite end: it buys the highest absolute
accuracy with the lowest efficiency, which is the expected and intended
trade-off between the two tiers.

\begin{table}[H]
\centering
\caption{Efficiency: balanced accuracy per million parameters (TTA BMA).}
\label{tab:efficiency}
\begin{tabular}{lrrrr}
\toprule
Model & Params (M) & GMACs & BMA & BMA / M-param \\
\midrule
ResNet-18            & 11.69 & 1.810 & 0.3368 & 0.029 \\
MobileNetV2          & 3.50  & 0.300 & 0.4738 & 0.135 \\
EfficientNet-B0      & 5.29  & 0.390 & 0.4946 & 0.093 \\
MobileViT-S          & 5.60  & 1.000 & 0.6020 & 0.108 \\
EfficientFormer-L1   & 12.27 & 1.300 & 0.5519 & 0.045 \\
Hybrid (ours, light) & 3.98  & 0.631 & 0.6081 & \textbf{0.153} \\
DEKAN (ours, flagship) & 16.45 & 6.633 & 0.6438 & 0.039 \\
\bottomrule
\end{tabular}
\end{table}

\begin{figure}[H]
\centering
\includegraphics[width=0.85\textwidth]{05_efficiency_scatter.pdf}
\caption{Accuracy vs.\ parameters; bubble size is proportional to GMACs. The
dotted line marks the 6M-parameter budget. The lightweight hybrid (red) gives the
best accuracy within budget.}
\label{fig:efficiency}
\end{figure}

\section{Effect of Test-Time Augmentation}

TTA helps most but not all models, so Table~\ref{tab:tta} reports BMA with and
without it. Both proposed models improve under TTA (the lightweight hybrid by
the largest margin among the strong models, $+0.027$). Two baselines are the
exception: EfficientFormer-L1 drops slightly, and ResNet-18 degrades sharply
($-0.151$). The ResNet-18 case is a known failure mode in which the per-view
predictions are confidently wrong in different directions, so averaging them
hurts; we report it openly rather than quietly dropping TTA for that model.

\begin{table}[H]
\centering
\caption{BMA without and with 8-view TTA, and the change.}
\label{tab:tta}
\begin{tabular}{lrrr}
\toprule
Model & BMA (no TTA) & BMA (TTA) & $\Delta$ \\
\midrule
ResNet-18            & 0.4875 & 0.3368 & $-0.151$ \\
MobileNetV2          & 0.4485 & 0.4738 & $+0.025$ \\
EfficientNet-B0      & 0.4917 & 0.4946 & $+0.003$ \\
MobileViT-S          & 0.5925 & 0.6020 & $+0.010$ \\
EfficientFormer-L1   & 0.5559 & 0.5519 & $-0.004$ \\
Hybrid (ours, light) & 0.5815 & 0.6081 & $+0.027$ \\
DEKAN (ours, flagship) & 0.6368 & 0.6438 & $+0.007$ \\
\bottomrule
\end{tabular}
\end{table}

\section{Analysis of Design Solutions (Ablations)}

\subsection{Lightweight hybrid}

Table~\ref{tab:ablhybrid} and Figure~\ref{fig:ablation} isolate the components of
the lightweight model. Going from the CNN stem alone (with global average pooling
and a linear head) to the full hybrid lifts BMA from 0.4853 to 0.6081, a gain of
0.123 --- a clear demonstration that the transformer trunk and metadata head
together carry most of the model's balanced accuracy. The effect of metadata in
isolation is, however, modest at test time on this single run: the full model
(0.6081) and the no-metadata variant (0.6085) are within 0.0004 of each other on
the test set, even though on the validation set the full model led by about
0.011 BMA. We therefore report the metadata contribution for the lightweight
model honestly as inconclusive on a single test evaluation; the DEKAN results
below give a clearer signal.

\begin{table}[H]
\centering
\caption{Lightweight hybrid ablation (TTA BMA). The pure-ViT (no CNN stem)
variant was not trained, so that row of the ablation is incomplete.}
\label{tab:ablhybrid}
\begin{tabular}{llr}
\toprule
Configuration & Variant & BMA (TTA) \\
\midrule
CNN stem + GAP + linear     & hybrid\_cnn\_only & 0.4853 \\
CNN + ViT, no metadata      & hybrid\_no\_meta  & 0.6085 \\
CNN + ViT + metadata (full) & hybrid\_full      & 0.6081 \\
Pure ViT, no CNN stem       & hybrid\_vit\_only & --- (not trained) \\
\bottomrule
\end{tabular}
\end{table}

\begin{figure}[H]
\centering
\includegraphics[width=0.85\textwidth]{06_ablation.pdf}
\caption{Lightweight hybrid ablation: BMA as components are added. (The pure-ViT
variant was not trained and is absent.)}
\label{fig:ablation}
\end{figure}

\subsection{DEKAN}

Table~\ref{tab:abldekan} and Figure~\ref{fig:dekanablation} ablate the flagship
model. Here metadata helps clearly: removing it drops BMA from 0.6438 to 0.6302,
a loss of 0.014. The KAN head, on the other hand, does not earn its place on this
task --- the otherwise identical model with a plain linear head actually scores
slightly higher (0.6510 vs.\ 0.6438). This is a genuine negative result that we
report as such: on ISIC-2019 at this scale, the KAN classifier did not improve
over a linear one.

\begin{table}[H]
\centering
\caption{DEKAN ablation (TTA BMA). The single-backbone variants were not trained,
so the dual-backbone fusion gain cannot yet be quantified.}
\label{tab:abldekan}
\begin{tabular}{llr}
\toprule
Configuration & Variant & BMA (TTA) \\
\midrule
Dual backbone + fusion, no metadata & dekan\_no\_meta       & 0.6302 \\
Dual backbone + fusion + metadata, linear head & dekan\_linear & 0.6510 \\
Dual backbone + fusion + metadata + KAN (full) & dekan\_full   & 0.6438 \\
DenseNet-only + metadata + KAN      & dekan\_densenet\_only & --- (not trained) \\
EfficientNet-only + metadata + KAN  & dekan\_effnet\_only   & --- (not trained) \\
\bottomrule
\end{tabular}
\end{table}

\begin{figure}[H]
\centering
\includegraphics[width=\textwidth]{07_dekan_ablation.pdf}
\caption{DEKAN ablation: effect of metadata and of the KAN head. (Single-backbone
variants were not trained and are absent.)}
\label{fig:dekanablation}
\end{figure}

\section{Comparisons and Relationships}

Read against the wider literature, these results sit where one would expect a
single-model, efficiency-constrained system to sit. The top entries on the
ISIC-2019 challenge leaderboard~\cite{isic2019leaderboard} reach a balanced
accuracy around 0.636 (DAISYLab) using large ensembles of multi-resolution
EfficientNets and SENet-154. DEKAN, a single model, reaches 0.6438 BMA on our
held-out split, and even the 3.98M-parameter lightweight hybrid reaches 0.6081 ---
within striking distance of an entire ensemble, at a tiny fraction of the cost.
The comparison is not exact, since the leaderboard is computed on the official
challenge test server and our split is an internal one, so we present it as
context rather than as a like-for-like ranking. Relative to controlled
single-model comparisons on this dataset such as Aktan et
al.~\cite{aktan2025comparison}, our contribution is to add clinical-metadata
fusion to the comparison and to report accuracy together with a hard efficiency
budget. The clearest relationship to prior work is with the metadata-fusion line
of Pacheco and Krohling~\cite{pacheco2020impact,pacheco2021attention}: we adopt
their core idea of representing missing values with learned embeddings, and we
extend their heavy-CNN setting to lightweight and dual-backbone hybrids.

\section{Discussions}

The results support the central claims of the thesis while leaving a few
questions open. The lightweight hybrid does what it was designed to do: it beats
all five same-budget baselines on balanced accuracy and is the most
parameter-efficient model in its class, which makes it a credible candidate for
deployment on modest hardware. DEKAN shows that the same design scales: investing
capacity in a second backbone, learned fusion, and a deeper trunk pushes balanced
accuracy to the level of a leaderboard ensemble in a single model. Metadata
fusion helps clearly in the higher-capacity model and is, at worst, harmless in
the lightweight one.

Two findings cut against initial expectations and are worth stating plainly.
First, the KAN classifier did not beat a linear head; the extra machinery of
learnable splines did not translate into accuracy on this task and scale.
Second, the benefit of metadata for the lightweight model could not be confirmed
on a single test evaluation, even though the validation curve favoured it. Both
are reported as they are rather than smoothed over.

\section{Limitations}

Several limitations bound these conclusions. All numbers come from a single seed
(42); publication-grade claims would require repeating each run over multiple
seeds and reporting mean and standard deviation, which we have not yet done.
Three planned runs were not completed --- the pure-ViT lightweight ablation and
the two single-backbone DEKAN ablations --- so the contribution of the CNN
inductive bias in isolation and the precise size of the dual-backbone fusion
gain remain unquantified. We also did not evaluate on an external dataset, so the
generalisation of these models beyond the ISIC-2019 distribution is untested.
Finally, the comparison with the ISIC-2019 leaderboard is contextual, not a
controlled head-to-head, because the evaluation splits differ.
